%% notation.tex
%% Symbol definitions and conventions for RMG Periodic RT-TDDFT paper (paper 2)
%% This file is \input{} into paper2_main.tex preamble area.
%% It defines macros AND documents the convention choices.
%% Never change a convention without updating DECISIONS.md first.

%% ============================================================
%% UNIT SYSTEM
%% ============================================================
%% Atomic units throughout (hbar=1, m_e=1, e=1, 4*pi*epsilon_0=1).
%% Stated explicitly at start of Section 2.
%% Exception: SI when quoting experimental values.

%% ============================================================
%% IMAGINARY UNIT
%% ============================================================
%% Use \imath (dotless i) consistent with paper 1.
%% Already provided by LaTeX. No new command needed.
%% Convention: paper 1 uses \imath throughout.

%% ============================================================
%% ELECTRON CHARGE SIGN CONVENTION
%% ============================================================
%% Electron charge = -e, where e > 0.
%% Current: \vb{J} = -e \vb{v} \rho  (explicit minus sign)
%% Stated at first use of J in text.

%% ============================================================
%% OPERATORS (hat notation, consistent with paper 1)
%% ============================================================

% Hamiltonian
\newcommand{\Hhat}{\hat{H}}

% Kohn-Sham Hamiltonian
\newcommand{\HKS}{\hat{H}_{\mathrm{KS}}}

% Time-evolution operator  U(t) = exp(Omega)
\newcommand{\Uop}{\mathcal{U}}

% Magnus superoperator
\newcommand{\Wop}{\hat{\mathcal{W}}}

% Momentum operator (canonical)  p = -i hbar grad
\newcommand{\phat}{\hat{\mathbf{p}}}

% Mechanical momentum  pi = p + (e/c) A
\newcommand{\pihat}{\hat{\boldsymbol{\pi}}}

% Velocity operator  v = pi / m_e  (includes [r, V^nl] correction)
\newcommand{\vhat}{\hat{\mathbf{v}}}

% Position operator
\newcommand{\rhat}{\hat{\mathbf{r}}}

% Current density operator
\newcommand{\Jhat}{\hat{\mathbf{J}}}

% Nonlocal pseudopotential operator
\newcommand{\Vnl}{V^{\mathrm{nl}}}

% Commutator shorthand  [A, B]
% Use standard \left[ \right] or physics package \comm{A}{B}
% physics package is loaded in paper2_main.tex

%% ============================================================
%% FIELDS AND POTENTIALS
%% ============================================================

% Vector potential
\newcommand{\Avec}{\mathbf{A}}

% Electric field
\newcommand{\Evec}{\mathbf{E}}

% Magnetic field
\newcommand{\Bvec}{\mathbf{B}}

% Scalar (electrostatic) potential
%\newcommand{\Phi}{\phi}

% Effective KS potential
\newcommand{\Veff}{V_{\mathrm{eff}}}

% Hartree potential
\newcommand{\VH}{V_{\mathrm{H}}}

% Exchange-correlation potential
\newcommand{\Vxc}{V_{\mathrm{xc}}}

% External (ionic) potential -- local part
\newcommand{\Vloc}{V^{\mathrm{loc}}}

%% ============================================================
%% DENSITY MATRIX AND PROPAGATION
%% ============================================================

% One-particle density matrix (consistent with paper 1: bold P)
% Paper 1 uses \vb{P} from physics package
% Use same convention here: \vb{P}, \vb{H}, etc. for matrices

% Density matrix at time t
%   P(t)  ->  \vb{P}(t)

% Initial density matrix
%   P^0 = P(0)  ->  \vb{P}^0

% Propagated density matrix
%   P(t+dt) = P^1  ->  \vb{P}^1

% Average Hamiltonian
%   Hbar = 1/2 (H0 + H1)  ->  \vb{\bar{H}}
\newcommand{\Hbar}{\bar{\mathbf{H}}}

% Magnus operator (first order)
%   Omega_1 = -i/hbar * Hbar * dt  ->  \Omega_1
% Use standard \Omega_1

% Curly bracket anti-commutator used in BCH expansion
%   {A,B} = AB + B^dag A  ->  defined in paper 1 notation
%   Keep same definition: \{A,B\} = AB + BA^\dag

%% ============================================================
%% BLOCH STATES AND K-POINTS
%% ============================================================

% Bloch wavefunction
%   psi_{n,k}(r,t)  ->  \psi_{n,\mathbf{k}}(\mathbf{r},t)

% Crystal momentum vector
%   k  ->  \mathbf{k}

% Lattice constant
%   L  (same as paper 1)

% Reciprocal lattice vector
%   G  ->  \mathbf{G}

% k-point weight
%   w_k  ->  w_{\mathbf{k}}

%% ============================================================
%% CURRENT AND OBSERVABLES
%% ============================================================

% Paramagnetic current density
\newcommand{\Jpara}{\mathbf{J}_{\mathrm{para}}}

% Diamagnetic current density
\newcommand{\Jdia}{\mathbf{J}_{\mathrm{dia}}}

% Total current density
\newcommand{\Jtot}{\mathbf{J}}

% Macroscopic current (volume-averaged, time-dependent observable)
%   J(t)  ->  \mathbf{J}(t)

% Optical conductivity tensor
%   sigma(omega)  ->  \boldsymbol{\sigma}(\omega)
\newcommand{\sigmaop}{\boldsymbol{\sigma}}

% Dielectric function
%   epsilon(omega)  ->  \boldsymbol{\varepsilon}(\omega)
\newcommand{\epsop}{\boldsymbol{\varepsilon}}

% Berry phase polarization
%   P_BP  ->  \mathbf{P}_{\mathrm{BP}}
\newcommand{\PBP}{\mathbf{P}_{\mathrm{BP}}}

%% ============================================================
%% MOMENTUM MATRIX ELEMENTS (current operator in KS basis)
%% ============================================================

% Momentum matrix element (paramagnetic part)
%   p^alpha_{ab} = <psi_a | (-i hbar grad_alpha + hbar k_alpha) | psi_b>
%   ->  p^{\alpha}_{ab}  (no special macro, use standard notation)

% Nonlocal PP correction to momentum matrix
%   [r_alpha, V^nl]_{ab} = <psi_a | [r_alpha, V^nl] | psi_b>

% Full current operator matrix (stored in code as Pxmatrix, Pymatrix, Pzmatrix)
%   J^alpha_{ab} = p^alpha_{ab} + [r_alpha, V^nl]_{ab}
\newcommand{\Jmat}[1]{\mathcal{J}^{#1}}
%   Usage: \Jmat{x}_{ab} for x-component matrix element

%% ============================================================
%% SPIN NOTATION
%% ============================================================

% Spin-up density matrix
\newcommand{\Pup}{P^{\uparrow}}

% Spin-down density matrix
\newcommand{\Pdn}{P^{\downarrow}}

% Spin-up electron density
\newcommand{\rhoup}{\rho^{\uparrow}}

% Spin-down electron density
\newcommand{\rhodn}{\rho^{\downarrow}}

% Spin magnetization density
\newcommand{\rhomag}{m_s}

%% ============================================================
%% PSEUDOPOTENTIAL PROJECTORS (Kleinman-Bylander)
%% ============================================================

% KB projector function at atom I, angular momentum lm
%   beta^{lm}_I(r)  ->  \beta^{lm}_{I}(\mathbf{r})

% KB projector overlap
%   <beta^{lm}_I | psi_n>  ->  standard bra-ket notation

%% ============================================================
%% GAUGE FUNCTION
%% ============================================================

% Gauge function (scalar field)
%   chi(r,t)  ->  \chi(\mathbf{r},t)
%   (standard \chi, no new command needed)

%% ============================================================
%% INHERITED FROM PAPER 1 (do not redefine)
%% ============================================================
%% These are already defined in paper 1 and must remain consistent:
%%   \vb{P}, \vb{H}  -- bold matrices (from physics package \vb{})
%%   \imath           -- imaginary unit
%%   \myred, \myblue, \mygrey, \mygreen  -- annotation colors
%%   \hbar            -- reduced Planck constant (standard LaTeX)
%%   All equation labels from paper 1 are cited as \cite{paper1}
%%   not cross-referenced (separate documents)
